\section{The uniform newborn growth exponent}
\label{uniform:section}

For $x\in\mathcal X$ put
\[
 F_x(t)=\mathbb E_xC_{\rm tag}(t),\qquad
 G(t)=\sup_xF_x(t),\qquad H(t)=1+\int_0^tG(r)\,dr.
\]
Every $F_x$ is nondecreasing, and $1\le F_x(t)\le e^t$. The same
bounds hold for $G$. At a deterministic or stopping time, partition
the future core genealogy into its currently existing cores and the
families issued from future children of these cores. Each existing
core creates children at intensity at most one, and each newborn
child starts in $\mathcal X$. The compensator formula gives
\begin{equation}\label{uniform:old-core}
 \mathbb E[C(t+u)\mid\mathcal F_t]\le C(t)H(u).
\end{equation}
This holds for any old retained descendant component in an admissible
environment, including the one-stub inert component. It does not
require independence between child families.

In particular, $G(t+u)\le G(t)H(u)$. Also $G(t)\le H(t)$ by the same
decomposition from one newborn core. Exchanging time parameters gives
$G(t+r)\le H(t)G(r)$, whence
\[
 H(t+u)\le H(t)+H(t)\int_0^uG(r)\,dr=H(t)H(u).
\]
Since $G$ is nondecreasing, $H(t)\le1+tG(t)$. These facts prove the
following standard submultiplicative consequence.

\begin{proposition}\label{uniform:lambda}
There is a finite exponent
\begin{equation}\label{uniform:lambda-equation}
 \lambda=\lim_{t\to\infty}\frac{\log G(t)}t
 =\lim_{t\to\infty}\frac{\log H(t)}t
 =\inf_{T>0}\frac{\log H(T)}T,
 \qquad \mu\le\lambda\le1.
\end{equation}
For each $\alpha>\lambda$, a finite constant $K_\alpha$ satisfies
$H(t)\le K_\alpha e^{\alpha t}$ for all $t\ge0$.
\end{proposition}
\begin{proof}
Apply continuous-time subadditivity to the locally bounded function
$\log H$, and use $G\le H\le1+tG$ to identify the two limits. The
inert newborn belongs to $\mathcal X$ and has exponent $\mu$ by
Theorem~\ref{gs:inert-mu}, proving the lower inequality. The upper
inequality follows from $H(t)\le e^t$. For the final assertion choose
$T$ with $H(T)<e^{\alpha T}$, write $t=nT+r$ with $0\le r<T$, and
use submultiplicativity and $H(r)\le e^T$.
\end{proof}

We will first use $\lambda$ without identifying it with $\mu$.
The equality is proved in Section~\ref{erc:closure-section}. This
order is essential: the controlled-moment estimates used in that
proof depend on $\lambda$, and are established independently of the
desired equality.
